% =========================================================
% COMPLEMENTO: INFERÊNCIA CAUSAL TEMPORAL
% Versão focada apenas no que é central para a narrativa:
% (1) macrofenótipos K=2 do artigo;
% (2) scan clínico B/M de benefício e malefício.
% =========================================================

\providecommand{\pp}{\,\text{pp}}
\providecommand{\h}{\,\text{h}}
\providecommand{\mcgkgmin}{\,\text{mcg/kg/min}}
\providecommand{\code}[1]{\texttt{#1}}

\section{Complemento metodológico: inferência causal temporal}

\subsection{Objetivo da etapa causal}

A etapa atual transforma o estudo anterior de fenotipagem em uma análise de inferência causal observacional. O artigo anterior identificou macrofenótipos fisiológicos a partir de trajetórias clínicas pré-transfusionais; agora, esses fenótipos são usados para perguntar se a associação entre transfusão de hemácias e mortalidade varia conforme o estado fisiológico do paciente antes da decisão transfusional.

O tratamento continua sendo a transfusão:

\[
A =
\begin{cases}
1, & \text{transfundido},\\
0, & \text{controle não transfundido}.
\end{cases}
\]

O desfecho primário é mortalidade. A pergunta causal é expressa em termos de desfechos potenciais:

\[
Y(1) = \text{desfecho caso o paciente fosse transfundido},
\]

\[
Y(0) = \text{desfecho caso o paciente não fosse transfundido}.
\]

Como a análise parte de pacientes transfundidos comparados a controles pareados, o estimando principal é o efeito médio do tratamento nos tratados:

\[
ATT_g = E\{Y(1)-Y(0)\mid A=1, G=g\},
\]

em que $G=g$ representa o macrofenótipo ou o subgrupo clínico identificado pelo scan.

\begin{boxsummary}
Em termos práticos, a análise estima, dentro de cada grupo fisiológico, se pacientes transfundidos tiveram mortalidade maior ou menor que controles pareados clinicamente comparáveis.
\end{boxsummary}

\subsection{Por que isso é inferência causal e não apenas clusterização}

A clusterização do artigo anterior foi usada para descobrir estrutura fisiológica latente. A inferência causal entra depois, quando comparamos transfundidos e controles dentro desses grupos.

A sequência conceitual é:

\[
\text{trajetória pré-}t_0
\rightarrow
\text{embedding fisiológico}
\rightarrow
\text{macrofenótipo / scan}
\rightarrow
\text{efeito da transfusão}.
\]

Assim, o cluster ou o grupo do scan não é o tratamento. O tratamento é a transfusão. O grupo fisiológico define a subpopulação na qual o efeito da transfusão é estimado.

\subsection{Ponto de decisão temporal}

O ponto de decisão $t_0$ foi definido como o instante da primeira transfusão nos pacientes tratados. Para os controles, foi utilizado um pseudo-$t_0$ compatível com a distribuição temporal dos tratados. A janela fisiológica principal considerada foi de 48 horas antes de $t_0$.

O desenho causal ideal exige que as variáveis usadas na estimação pertençam ao período pré-intervenção:

\[
X_{\text{pre}} = \{X_t: t < t_0\}.
\]

Portanto, a estrutura temporal desejada é:

\[
X_{\text{pre}} \longrightarrow A \longrightarrow Y.
\]

\begin{boxsummary}
Nota de cautela: a execução reaproveitada indicava \code{keep\_post\_t0\_features=true}. Por isso, os resultados atuais devem ser apresentados como exploração causal baseada nos artefatos do artigo. A versão confirmatória deve remover estritamente qualquer variável derivada após $t_0$.
\end{boxsummary}

\subsection{Hipóteses necessárias para interpretação causal}

Para interpretar a diferença pareada como efeito causal, são necessárias hipóteses fortes:

\begin{itemize}[leftmargin=1.5em, itemsep=0.25em]
  \item \textbf{Consistência:} o desfecho observado corresponde ao desfecho potencial sob o tratamento recebido.
  \item \textbf{Ignorabilidade condicional:} após ajustar por covariáveis, gravidade e trajetória fisiológica, não há confundidor não medido que determine simultaneamente transfusão e mortalidade.
  \item \textbf{Positividade:} dentro de cada grupo, existem pacientes comparáveis transfundidos e não transfundidos.
  \item \textbf{Ausência de vazamento temporal:} informações após $t_0$ não devem entrar nas variáveis usadas para estimar o efeito.
  \item \textbf{Comparabilidade clínica:} o pareamento deve aproximar tratados e controles em fatores de gravidade e indicação transfusional.
\end{itemize}

Em UTI, a hipótese mais frágil é a ignorabilidade condicional, pois a decisão transfusional pode depender de sangramento ativo, julgamento clínico, instabilidade aguda, procedimentos e metas específicas de hemoglobina não totalmente medidas. Por isso, os resultados devem ser descritos como \textbf{inferência causal observacional exploratória}.

\section{População e resultado global}

A coorte importada do pipeline de fenotipagem continha 40.167 internações em UTI, sendo 10.639 transfundidas e 29.528 controles. A estimativa pareada principal utilizou 10.009 pares tratados-controles.

\begin{table}[ht]
\centering
\caption{Resumo da população analisada.}
\begin{tabular}{lr}
\toprule
Componente & Valor \\
\midrule
Internações totais & 40.167 \\
Transfundidos & 10.639 \\
Controles & 29.528 \\
Pares pareados & 10.009 \\
Janela temporal & 48 h \\
Representação temporal & MiniRocket \\
Desfecho primário & Mortalidade \\
\bottomrule
\end{tabular}
\end{table}

Na população pareada completa, a mortalidade foi menor nos transfundidos:

\[
\widehat{ATT}_{global}
= 0{,}385 - 0{,}446
= -0{,}061.
\]

Isso corresponde a uma diferença de aproximadamente $-6{,}1\pp$.

\begin{table}[ht]
\centering
\caption{Estimativa global pareada.}
\begin{tabular}{lrrrrr}
\toprule
Desfecho & Pares & Transfundidos & Controles & Diferença & IC95\% \\
\midrule
Mortalidade & 10.009 & 0,385 & 0,446 & -0,061 & [-0,075;\,-0,047] \\
Ventilação & 10.009 & 64,3 h & 32,9 h & +31,4 h & [+29,0;\,+33,9] \\
RRT & 10.009 & 0,104 & 0,065 & +0,039 & [+0,032;\,+0,047] \\
Noradrenalina eq. máx. & 10.009 & 0,265 & 0,200 & +0,065 & [+0,031;\,+0,111] \\
Permanência UTI & 10.009 & 732,8 h & 421,4 h & +311,4 h & [+291,9;\,+330,9] \\
\bottomrule
\end{tabular}
\end{table}

O resultado global sugere menor mortalidade associada à transfusão, mas também maior carga de suporte orgânico. Portanto, o efeito global não deve ser interpretado como homogêneo. Ele mistura pacientes potencialmente beneficiados, pacientes neutros e pacientes em que a transfusão pode ser marcador de gravidade ou estar associada a pior prognóstico.

\section{Camada principal do artigo: macrofenótipos K=2}

A solução $K=2$ é a camada central herdada do artigo anterior. Ela resume os pacientes em dois macrofenótipos fisiológicos: um grupo com benefício relativo e outro com maior carga de suporte orgânico.

\begin{table}[ht]
\centering
\caption{Efeito sobre mortalidade nos macrofenótipos K=2.}
\begin{tabular}{lrrrrr}
\toprule
Macrofenótipo & Pares & Transfundidos & Controles & Diferença & IC95\% \\
\midrule
Benefício relativo & 7.085 & 0,357 & 0,423 & -0,066 & [-0,081;\,-0,049] \\
Alto suporte / risco & 2.924 & 0,451 & 0,501 & -0,050 & [-0,077;\,-0,024] \\
\bottomrule
\end{tabular}
\end{table}

A interpretação do $K=2$ é importante porque ele organiza a narrativa fisiológica do artigo. O primeiro macrofenótipo concentra pacientes com menor carga de suporte e sinal de benefício relativo. O segundo macrofenótipo não apresentou aumento de mortalidade nos transfundidos, mas apresentou grande aumento de suporte orgânico: aproximadamente $+73{,}3\h$ de ventilação, $+9{,}6\pp$ de RRT, $+0{,}173\mcgkgmin$ de noradrenalina equivalente máxima e $+776{,}2\h$ de permanência em UTI.

\begin{boxsummary}
Assim, $K=2$ deve ser apresentado como macrofenotipagem clínica do artigo: ele mostra que existe heterogeneidade fisiológica relevante, mas não é a camada mais fina para nomear benefício ou malefício.
\end{boxsummary}

\section{Camada refinada: scan de benefício e malefício}

O scan foi usado para refinar a interpretação dos macrofenótipos. Em vez de discutir escolhas alternativas de clusterização, a análise relevante aqui é o resultado clínico do scan: identificação de subgrupos com sinal favorável ou desfavorável de mortalidade.

O critério básico foi:

\[
\text{diff\_mort}
=
\text{Mort}_{\text{transfundido}}
-
\text{Mort}_{\text{controle}}.
\]

\begin{itemize}[leftmargin=1.5em, itemsep=0.25em]
  \item $\text{diff\_mort}<0$: sinal de benefício associativo.
  \item $\text{diff\_mort}>0$: sinal de malefício ou risco associativo.
\end{itemize}

\subsection{Subgrupos com sinal de benefício}

\begin{table}[ht]
\centering
\caption{Subgrupos do scan com sinal favorável.}
\begin{tabular}{llrrr}
\toprule
Grupo & Leitura clínica resumida & $n$ total & $n$ repr. & Diferença repr. \\
\midrule
B1 & Perfundido amplo / baixa taquicardia & 5.476 & 1.401 & -22,48 pp \\
B2 & Anemia dinâmica / oxigenação estável & 2.585 & 1.033 & -33,59 pp \\
B3 & Estresse metabólico com perfusão preservada & 1.921 & 921 & -32,46 pp \\
B4 & Pressão média elevada / contexto neurovascular & 1.369 & 979 & -31,97 pp \\
\bottomrule
\end{tabular}
\end{table}

Os grupos B compartilham uma hipótese fisiológica: a transfusão parece ocorrer em pacientes nos quais a oferta de oxigênio pode estar limitada, mas sem assinatura dominante de falência orgânica irreversível. A leitura clínica mais plausível é de pacientes \textbf{DO$_2$-limitados}, com anemia corrigível, oxigenação relativamente preservada ou contexto vascular/neurovascular sensível à anemia.

\begin{boxbenefit}{Interpretação causal dos grupos B}
Os grupos B não provam benefício causal definitivo, mas concentram subpopulações nas quais a diferença pareada favorece os transfundidos. A hipótese é que, nesses pacientes, a transfusão esteja mais alinhada à correção de oferta de oxigênio do que à resposta tardia a falência orgânica avançada.
\end{boxbenefit}

\subsection{Subgrupos com sinal de malefício}

\begin{table}[ht]
\centering
\caption{Subgrupos do scan com sinal desfavorável.}
\begin{tabular}{llrrr}
\toprule
Grupo & Leitura clínica resumida & $n$ total & $n$ repr. & Diferença repr. \\
\midrule
M1 & Falência hepatobiliar / estresse metabólico & 654 & 654 & +26,61 pp \\
M2 & Doença crítica sustentada hepatobiliar-renal & 1.129 & 1.129 & +25,86 pp \\
M3 & Baixo fluxo cardiorrenal / DBP baixa & 913 & 913 & +24,32 pp \\
M4 & Falência progressiva renal-metabólica & 753 & 753 & +23,24 pp \\
\bottomrule
\end{tabular}
\end{table}

Os grupos M compartilham assinatura de gravidade avançada: disfunção hepatobiliar, disfunção renal, baixo fluxo, lactato em piora, necessidade de suporte e maior probabilidade de transfusão como marcador de trajetória crítica.

\begin{boxharm}{Interpretação causal dos grupos M}
Os grupos M indicam maior mortalidade ou pior prognóstico entre transfundidos pareados. A interpretação principal é confusão por indicação e marcador de gravidade, mas existe plausibilidade de dano em subcenários de congestão, sobrecarga, inflamação ou falência orgânica avançada.
\end{boxharm}

\section{Como K=2 e scan se complementam}

A leitura final deve combinar as duas camadas:

\begin{itemize}[leftmargin=1.5em, itemsep=0.25em]
  \item \textbf{K=2} fornece a macroestrutura fisiológica do artigo: pacientes mais compensados versus pacientes de maior suporte e complexidade.
  \item \textbf{Scan} fornece a tradução clínica fina: grupos B com sinal favorável e grupos M com sinal desfavorável.
\end{itemize}

Em outras palavras, $K=2$ mostra que há heterogeneidade fisiológica relevante; o scan mostra quais padrões específicos dentro dessa heterogeneidade parecem concentrar benefício ou risco.

\begin{boxsummary}
A conclusão central é que a transfusão de hemácias não deve ser analisada como intervenção homogênea. O sinal associado à transfusão depende do estado fisiológico pré-intervenção: perfis DO$_2$-limitados parecem concentrar sinal favorável, enquanto perfis de falência orgânica avançada ou baixo fluxo cardiorrenal concentram sinal desfavorável.
\end{boxsummary}

\section{Modelos contrafactuais individuais}

Além do pareamento e do scan, foi implementado um modelo contrafactual inicial com dois modelos preditivos:

\[
\hat{Y}_i(1)=f_1(X_i),
\qquad
\hat{Y}_i(0)=f_0(X_i).
\]

O efeito individual estimado foi:

\[
\widehat{ITE}_i = \hat{Y}_i(1)-\hat{Y}_i(0).
\]

Esse modelo estima, para cada paciente, o risco esperado sob transfusão e sob não transfusão. O desempenho preditivo para mortalidade foi:

\begin{table}[ht]
\centering
\caption{Desempenho preditivo do modelo contrafactual inicial.}
\begin{tabular}{lrr}
\toprule
Métrica & Completo & Teste \\
\midrule
AUC & 0,869 & 0,849 \\
Average precision & 0,780 & 0,732 \\
Brier score & 0,135 & 0,145 \\
Acurácia & 0,808 & 0,789 \\
F1 & 0,673 & 0,637 \\
\bottomrule
\end{tabular}
\end{table}

Esse resultado mostra que as covariáveis fisiológicas têm capacidade preditiva razoável para mortalidade. Entretanto, boa predição não garante validade causal do contrafactual. O modelo deve ser usado como análise exploratória e geração de hipóteses individuais, não como regra clínica.

\section{Limitações}

\begin{enumerate}[leftmargin=1.5em, itemsep=0.3em]
  \item \textbf{Confusão por indicação:} pacientes transfundidos podem ter sangramento, instabilidade, procedimento ou julgamento clínico não completamente capturado.
  \item \textbf{Vazamento temporal potencial:} a execução reaproveitada manteve algumas features pós-$t_0$; a análise confirmatória deve usar apenas variáveis pré-$t_0$.
  \item \textbf{Scan exploratório:} os grupos B/M foram descobertos nos próprios dados e precisam de validação.
  \item \textbf{Overlap:} alguns subgrupos podem ter controles comparáveis limitados.
  \item \textbf{Desfechos de suporte:} maior ventilação ou permanência pode refletir dano, maior sobrevida ou gravidade residual.
  \item \textbf{Contrafactual não observado:} $Y(1)$ e $Y(0)$ nunca são observados simultaneamente no mesmo paciente.
\end{enumerate}

\section{Próxima análise confirmatória}

A etapa confirmatória deve manter a narrativa focada em $K=2$ e scan, mas reforçar a validade causal:

\begin{enumerate}[leftmargin=1.5em, itemsep=0.25em]
  \item reconstruir features estritamente pré-$t_0$;
  \item recalcular propensity score e pareamento dentro dos macrofenótipos e dos grupos B/M;
  \item reportar SMD antes e depois do ajuste;
  \item avaliar overlap por grupo;
  \item repetir bootstrap para os efeitos B/M;
  \item testar sensibilidade a caliper e à janela temporal;
  \item verificar se os grupos B/M mantêm direção em split temporal ou coorte externa.
\end{enumerate}

\section{Frase-síntese para discussão}

\begin{quote}
Nesta análise observacional pareada, a transfusão de hemácias apresentou associação heterogênea com mortalidade. O macrofenótipo $K=2$ confirmou a existência de dois estados fisiológicos amplos, enquanto o scan refinou essa heterogeneidade em subgrupos com sinal favorável e desfavorável. Os grupos B sugerem perfis potencialmente DO$_2$-limitados, nos quais a transfusão pode estar associada a menor mortalidade; os grupos M sugerem perfis de falência orgânica avançada ou baixo fluxo cardiorrenal, nos quais a transfusão pode refletir gravidade extrema e se associar a pior prognóstico.
\end{quote}
